\documentclass[11pt,a4paper]{article}

\usepackage[T1]{fontenc}
\usepackage[utf8]{inputenc}
\usepackage{lmodern}
\usepackage{microtype}
\usepackage{amsmath,amssymb,amsthm,mathtools}
\usepackage{geometry}
\usepackage{enumitem}
\usepackage{booktabs}
\usepackage{tikz-cd}
\usepackage[colorlinks=true,linkcolor=blue,citecolor=blue,urlcolor=blue]{hyperref}

\geometry{margin=2.65cm}
\setlength{\emergencystretch}{3em}
\setlist[itemize]{leftmargin=2.2em}
\setlist[enumerate]{leftmargin=2.2em}

\newtheorem{theorem}{Theorem}[section]
\newtheorem{proposition}[theorem]{Proposition}
\newtheorem{lemma}[theorem]{Lemma}
\newtheorem{corollary}[theorem]{Corollary}
\theoremstyle{definition}
\newtheorem{definition}[theorem]{Definition}
\newtheorem{example}[theorem]{Example}
\theoremstyle{remark}
\newtheorem{remark}[theorem]{Remark}

\newcommand{\OO}{\mathcal O}
\newcommand{\II}{\mathcal I}
\newcommand{\FF}{\mathcal F}
\newcommand{\GG}{\mathcal G}
\newcommand{\EE}{\mathcal E}
\newcommand{\PP}{\mathbf P}
\newcommand{\Spec}{\operatorname{Spec}}
\newcommand{\Coh}{\operatorname{Coh}}
\newcommand{\QCoh}{\operatorname{QCoh}}
\newcommand{\Hom}{\operatorname{Hom}}
\newcommand{\Ext}{\operatorname{Ext}}
\newcommand{\shHom}{\mathcal H\!om}
\newcommand{\shExt}{\mathcal E\!xt}
\newcommand{\RHom}{\mathbf R\!\operatorname{Hom}}
\newcommand{\RshHom}{\mathbf R\!\mathcal H\!om}
\newcommand{\RGamma}{\mathbf R\Gamma}
\newcommand{\D}{\mathrm D}
\newcommand{\Dqc}{\mathrm D_{\mathrm{qc}}}
\newcommand{\Dbcoh}{\mathrm D^b_{\mathrm{coh}}}
\newcommand{\depth}{\operatorname{depth}}
\newcommand{\grade}{\operatorname{grade}}
\newcommand{\height}{\operatorname{ht}}
\newcommand{\pd}{\operatorname{pd}}
\newcommand{\tr}{\operatorname{tr}}
\newcommand{\codim}{\operatorname{codim}}
\newcommand{\Supp}{\operatorname{Supp}}
\newcommand{\dual}{^{\vee}}
\newcommand{\determinant}{\operatorname{det}}

\title{Serre Duality for Cohen--Macaulay Schemes\\
       and Local Complete Intersections}
\author{Andrea Taccani}
\date{Seminar on Cohomology of Sheaves and Schemes, SoSe 2026}

\begin{document}
\maketitle

\begin{abstract}
Let $X$ be a projective, equidimensional Cohen--Macaulay scheme of dimension
$d$ over a field $k$, and let $\omega_X$ be a dualizing sheaf. The purpose of
this report is to explain why the degree-$d$ representing property of
$\omega_X$ extends functorially to all Ext groups:
\[
  \Ext_X^i(\FF,\omega_X)\simeq H^{d-i}(X,\FF)\dual
\]
for every coherent sheaf $\FF$ and every $i\geq0$, with the convention that
$H^q(X,\FF)=0$ for $q<0$. Choose a projective embedding
$i:X\hookrightarrow\PP^N_k$. The proof is organized through its dualizing complex
\[
  i_*\omega_X^\bullet
  \simeq\RshHom_{\PP^N_k}(i_*\OO_X,\omega_{\PP^N_k})[N].
\]
Projective-space duality gives the derived adjunction, while the Cohen--Macaulay condition implies the concentration $\omega_X^\bullet\simeq\omega_X[d]$. For a local complete intersection the Koszul resolution computes the unique cohomology sheaf and yields
\[
  \omega_X\simeq
  \omega_{\PP^N_k}|_X\otimes\determinant(\II/\II^2)\dual.
\]
The smooth formula $\omega_X\simeq\bigwedge^d\Omega_{X/k}$ follows by taking determinants in the conormal sequence.
\end{abstract}

\tableofcontents

\section{From the dualizing sheaf to full Serre duality}

Let $k$ be a field, let $X$ be a projective scheme of pure dimension $d$, and write
\[
  f:X\longrightarrow \Spec k
\]
for the structure morphism. A dualizing sheaf on $X$ is a coherent sheaf $\omega_X$ together with functorial isomorphisms
\begin{equation}\label{eq:dualizing-sheaf-definition}
  \Hom_X(\FF,\omega_X)\xrightarrow{\ \sim\ }H^d(X,\FF)\dual
\end{equation}
for all coherent sheaves $\FF$. Equivalently, it represents the contravariant functor $\FF\mapsto H^d(X,\FF)\dual$. Every projective scheme over a field admits such a sheaf; see
\cite[Chapter III, Proposition 7.5]{HartshorneAG} and
\cite[Tag 0FVV]{Stacks}.

The problem is therefore not the existence of a degree-$d$ representing
object. It is to determine when the same sheaf controls all higher Ext
groups. For a projective Cohen--Macaulay scheme the desired statement, for a
coherent sheaf $\FF$ and $i\geq0$, is
\begin{equation}\label{eq:target-duality}
  \Ext_X^i(\FF,\omega_X)\simeq H^{d-i}(X,\FF)\dual.
\end{equation}
Here and below, negative-degree sheaf cohomology is understood to be zero.

The precise statement, including the description of the perfect pairing, is
given and proved in Section~\ref{sec:serre-duality-pairing}.

\section{Derived categories and projective-space duality}

The construction of Ext, sheaf Ext, and their long exact sequences will be
used without repetition; see
\cite[Chapter III, Section 6]{HartshorneAG},
\cite[Sections 10.4 and 10.7]{Weibel}, and
\cite[Tags 06XP and 08DH]{Stacks}.

\begin{definition}[Complexes and derived categories]
\label{def:derived-categories}
Let $Y$ be a noetherian scheme, and let $\operatorname{Mod}(\OO_Y)$ denote
the abelian category of sheaves of $\OO_Y$-modules. A \emph{cochain complex}
$K^\bullet$ in this category is a sequence
\[
  \cdots\longrightarrow K^{n-1}\xrightarrow{d^{n-1}}K^n
  \xrightarrow{d^n}K^{n+1}\longrightarrow\cdots,
  \qquad d^n d^{n-1}=0.
\]
Its $n$-th cohomology sheaf is
$H^n(K^\bullet)=\ker(d^n)/\operatorname{im}(d^{n-1})$. A morphism of
complexes $f:K^\bullet\to L^\bullet$ is a family of maps
$f^n:K^n\to L^n$ for which, for every $n$, $ f^{n+1}\circ d_K^n=d_L^n\circ f^n$.

Such a morphism is a \emph{quasi-isomorphism} if each induced map
$H^n(f):H^n(K^\bullet)\to H^n(L^\bullet)$ is an isomorphism. Write
$\mathrm K(\OO_Y)$ for the homotopy category: its objects are
complexes and its morphisms are morphisms of complexes modulo chain homotopy.
Explicitly, two morphisms $f,g:K^\bullet\to L^\bullet$ are homotopic if there are maps
$h^n:K^n\to L^{n-1}$ such that
\[
  f^n-g^n=d_L^{n-1}h^n+h^{n+1}d_K^n
  \qquad\text{for every }n.
\]
If $Q$ is the class of quasi-isomorphisms, the derived category is the localization
\[
  \D(\OO_Y):=\mathrm K(\OO_Y)[Q^{-1}].
\]
In particular, a quasi-isomorphism becomes an isomorphism in $\D(\OO_Y)$.

The full subcategories $\D^+(\OO_Y)$ and $\D^b(\OO_Y)$ consist respectively
of those $K\in\D(\OO_Y)$ for which $H^n(K)=0$ for $n\ll0$, and those for
which $H^n(K)=0$ for all but finitely many $n$. Then
\[
  \D^+_{\mathrm{qc}}(Y)
  :=\{K\in\D^+(\OO_Y):H^n(K)\text{ is quasi-coherent for every }n\}
\]
and
\[
  \Dbcoh(Y)
  :=\{K\in\D^b(\OO_Y):H^n(K)\text{ is coherent for every }n\}.
\]
These are conditions on the cohomology sheaves, not on the terms of a chosen
representative. Since $Y$ is noetherian, a quasi-coherent sheaf is coherent
if and only if its module on every affine open is finitely generated. See
\cite[Tag 01XY]{Stacks}.

We use the cohomological shift convention
\[
  K[m]^r=K^{m+r},\qquad
  d_{K[m]}^r=(-1)^m d_K^{m+r}.
\]
Thus a sheaf $M$ placed in degree zero has its only nonzero cohomology sheaf
in $M[m]$ in degree $-m$.
\end{definition}

\begin{definition}[Internal and global derived Hom]
\label{def:derived-hom}
For complexes $K^\bullet,L^\bullet$ of $\OO_Y$-modules, their internal Hom
complex is defined by
\[
  \shHom_Y^n(K^\bullet,L^\bullet)
  :=\prod_{p\in\mathbf Z}\shHom_{\OO_Y}(K^p,L^{p+n}),
\]
with differential
\[
  d(\varphi)=d_L\circ\varphi-(-1)^n\varphi\circ d_K
  \qquad(\varphi\text{ of degree }n).
\]
For $K,L\in\D(\OO_Y)$, first choose complexes $K^\bullet$ and $L^\bullet$
representing them. A \emph{K-injective resolution} of $L^\bullet$ consists
of a complex $I^\bullet$ and a morphism of complexes
\[
  \iota:L^\bullet\longrightarrow I^\bullet
\]
satisfying two separate conditions. First, $\iota$ is a quasi-isomorphism, thus it becomes an isomorphism in $\D(\OO_Y)$, so $I^\bullet$ is another representative of the object $L$.
Second, $I^\bullet$ is K-injective, meaning that
\[
  \Hom_Y^\bullet(A^\bullet,I^\bullet)
  :=\Gamma\bigl(Y,\shHom_Y^\bullet(A^\bullet,I^\bullet)\bigr)
\]
is acyclic for every acyclic complex $A^\bullet$. This second condition is
what permits derived Hom into $L$ to be computed using ordinary Hom into
$I^\bullet$. The \emph{internal derived Hom} is therefore
\[
  \RshHom_Y(K,L):=\shHom_Y^\bullet(K^\bullet,I^\bullet)
  \quad\text{in }\D(\OO_Y).
\]
This object is independent, up to a unique functorial isomorphism, of the
chosen representatives. Equivalently, it is characterized by the natural
derived tensor--Hom adjunction
\[
  \Hom_{\D(\OO_Y)}
    \bigl(A,\RshHom_Y(K,L)\bigr)
  \simeq
  \Hom_{\D(\OO_Y)}
    \bigl(A\otimes_{\OO_Y}^{\mathbf L}K,L\bigr).
\]

The \emph{global derived Hom} is the derived global-sections complex
\[
  \RHom_Y(K,L)
  :=\RGamma\bigl(Y,\RshHom_Y(K,L)\bigr).
\]
When $Y$ is a $k$-scheme, as throughout this report, it is naturally an
object of $\D(k)$. If $K^\bullet$ is chosen K-flat, meaning that tensoring it
with any acyclic complex produces an acyclic complex, and $I^\bullet$ is
K-injective, then $\shHom_Y^\bullet(K^\bullet,I^\bullet)$ is K-injective.
Consequently, ordinary global sections compute its derived global sections,
and the global derived Hom is represented by the ordinary Hom complex
\[
  \Hom_Y^\bullet(K^\bullet,I^\bullet)
  =\Gamma\bigl(Y,\shHom_Y^\bullet(K^\bullet,I^\bullet)\bigr).
\]
For every $n\in\mathbf Z$ there are canonical identifications
\[
  H^n\RshHom_Y(K,L)=\shExt_Y^n(K,L),
  \qquad
  H^n\RHom_Y(K,L)
  =\Hom_{\D(\OO_Y)}(K,L[n]).
\]
See \cite[Tags 08DH and 0B6A]{Stacks}. By contrast,
$\Hom_Y(\FF,\GG)$ denotes ordinary morphisms when both arguments are sheaves.
\end{definition}

\begin{proposition}[Ext and cohomology as derived morphisms]
\label{prop:ext-derived}
For coherent sheaves $\FF,\GG$ on a noetherian $k$-scheme $Y$ and $i\geq0$,
there are canonical isomorphisms
\[
  \Ext_Y^i(\FF,\GG)
  \simeq \Hom_{\D^+_{\mathrm{qc}}(Y)}(\FF,\GG[i])
\]
and
\[
  H^i(Y,\FF)
  \simeq \Hom_{\D^+_{\mathrm{qc}}(Y)}(\OO_Y,\FF[i]).
\]
\end{proposition}

\begin{proof}
By the definition of Ext in a derived category,
\[
  \Ext_Y^i(\FF,\GG)=\Hom_{\D(\OO_Y)}(\FF,\GG[i]);
\]
see \cite[Section 10.7]{Weibel} and \cite[Tags 06XQ and 06XP]{Stacks}.
Both objects lie in the full subcategory $\D^+_{\mathrm{qc}}(Y)$, so the Hom
group may be computed there. For the second assertion, the equality of left
exact functors
$\Gamma(Y,-)=\Hom_{\OO_Y}(\OO_Y,-)$ gives a canonical isomorphism of their
right-derived functors
\[
  \RGamma(Y,-)
  \xrightarrow{\ \sim\ }
  \RHom_Y(\OO_Y,-)
  \qquad\text{in }\D(k).
\]

Taking degree-$i$ cohomology and using Definition \ref{def:derived-hom} gives
\begin{align*}
  H^i(Y,\FF)
  &=H^i\RGamma(Y,\FF)\\
  &\simeq H^i\RHom_Y(\OO_Y,\FF)\\
  &=\Hom_{\D(\OO_Y)}(\OO_Y,\FF[i])\\
  &=\Hom_{\D^+_{\mathrm{qc}}(Y)}(\OO_Y,\FF[i]).
\end{align*}
\end{proof}

\begin{lemma}[Duality over a field]
\label{lem:field-duality}
Let $C\in\D^b(k)$ have finite-dimensional cohomology. For every
$j\in\mathbf Z$ there is a natural isomorphism
\[
  \Hom_{\D(k)}(C,k[j])\simeq H^{-j}(C)\dual.
\]
\end{lemma}

\begin{proof}
Since $k$ is a field, the functor $\Hom_k(-,k)$ is exact. Hence
$\RHom_k(C,k)$ is represented by the dual complex
\[
  C^\vee:=\Hom_k^\bullet(C,k),
  \qquad
  (C^\vee)^j=\Hom_k(C^{-j},k).
\]
Exactness also implies that dualization commutes with taking cohomology.
Therefore
\begin{align*}
  \Hom_{\D(k)}(C,k[j])
    &=H^0\!\left(\RHom_k(C,k[j])\right)\\
    &\simeq H^j\!\left(\RHom_k(C,k)\right)\\
    &=H^j(C^\vee)\\
    &\simeq\Hom_k\!\left(H^{-j}(C),k\right)
     =H^{-j}(C)^\vee.
\end{align*}
Every map in this calculation is canonical and compatible with morphisms
$C\to C'$; hence the resulting isomorphism is natural in $C$. See
\cite[Sections 10.4 and 10.7]{Weibel}.
\end{proof}

\begin{lemma}[Extension from sheaves to bounded coherent complexes]
\label{lem:devissage}
Let $Y$ be noetherian, and let $(T^j)_{j\in\mathbf Z}$ and
$(S^j)_{j\in\mathbf Z}$ be contravariant cohomological functors on
$\Dbcoh(Y)$. Suppose
\[
  \nu^j:T^j\Longrightarrow S^j,
  \qquad j\in\mathbf Z,
\]
is a morphism of cohomological functors: each $\nu^j$ is natural, and the
maps $\nu^j$ commute with the connecting morphisms in the long exact
sequences associated with distinguished triangles. If
$\nu^j_{\mathcal F}$ is an isomorphism for every coherent sheaf
$\mathcal F$ placed in degree zero and every $j\in\mathbf Z$, then
$\nu^j_G$ is an isomorphism for every $G\in\Dbcoh(Y)$ and every
$j\in\mathbf Z$.
\end{lemma}

\begin{proof}
Proceed by induction on the number of nonzero cohomology sheaves of $G$.
Choose the largest integer $n$ for which $H^n(G)\neq0$. The standard
truncation triangle
\[
  \tau_{\leq n-1}G\longrightarrow G\longrightarrow H^n(G)[-n]
  \longrightarrow\tau_{\leq n-1}G[1]
\]
expresses $G$ in terms of a complex with one fewer nonzero cohomology sheaf
and a shifted coherent sheaf. Compatibility of $\nu^\bullet$ with the
connecting morphisms gives a morphism between the two associated long exact
sequences. By induction, the vertical maps for $\tau_{\leq n-1}G$ are
isomorphisms. The maps for $H^n(G)[-n]$ are also isomorphisms, because a
shift only reindexes the cohomological degree and $H^n(G)$ is coherent. The
five lemma now shows that $\nu^j_G$ is an isomorphism for every $j$.
\end{proof}

\begin{theorem}[Projective-space duality in derived form]
\label{thm:projective-derived}
Let $p:P=\PP^N_k\to\Spec k$. For every $G\in\Dbcoh(P)$ and every
$j\in\mathbf Z$, there is a functorial isomorphism
\begin{equation}\label{eq:projective-derived}
  \Hom_{\Dbcoh(P)}\bigl(G,\omega_P[N+j]\bigr)
  \simeq
  \Hom_{\D(k)}\bigl(\RGamma(P,G),k[j]\bigr).
\end{equation}
\end{theorem}

\begin{proof}
For $G\in\Dbcoh(P)$, set
\[
  T^j(G):=\Hom_{\Dbcoh(P)}\bigl(G,\omega_P[N+j]\bigr),
  \qquad
  S^j(G):=\Hom_{\D(k)}\bigl(\RGamma(P,G),k[j]\bigr).
\]
Projective-space duality supplies a trace morphism
\[
  \operatorname{Tr}_P:
  \RGamma\bigl(P,\omega_P[N]\bigr)\longrightarrow k
\]
whose map on degree-zero cohomology is the classical trace
$H^N(P,\omega_P)\to k$. For
\[
  \alpha:G\longrightarrow\omega_P[N+j],
\]
apply $\RGamma(P,-)$ and then compose with the shifted trace:
\[
  \RGamma(P,G)
  \xrightarrow{\RGamma(\alpha)}
  \RGamma\bigl(P,\omega_P[N+j]\bigr)
  \simeq
  \RGamma\bigl(P,\omega_P[N]\bigr)[j]
  \xrightarrow{\operatorname{Tr}_P[j]}
  k[j].
\]
Define
\[
  \eta_G^j(\alpha)
  :=\operatorname{Tr}_P[j]\circ\RGamma(\alpha).
\]
This construction is natural in $G$. Moreover, because $\RGamma(P,-)$ is a
triangulated functor and composition with the fixed morphism
$\operatorname{Tr}_P[j]$ respects shifts, the family
\[
  \eta^j:T^j\Longrightarrow S^j
\]
is compatible with the connecting morphisms; it is therefore a morphism of
cohomological functors in the sense of Lemma \ref{lem:devissage}. See
\cite[Tag 071J]{Stacks} for the fact that derived global sections is an exact
functor of derived categories, and \cite[Tags 014V and 0149]{Stacks} for
compatibility with shifts and the cohomologicality of representable Hom
functors.

Suppose that $G=\mathcal G$ is a coherent sheaf placed in degree zero. If
$j\geq-N$, then $N+j\geq0$, and $\eta_{\mathcal G}^j$ is precisely the
classical Serre-duality morphism
\[
  \Ext_P^{N+j}(\mathcal G,\omega_P)
  \longrightarrow H^{-j}(P,\mathcal G)^\vee.
\]
Under Lemma~\ref{lem:field-duality}, its target is
$S^j(\mathcal G)$, and classical projective-space duality says that it is an
isomorphism; see \cite[Chapter III, Theorem 7.1]{HartshorneAG}.

If $j<-N$, then $N+j<0$, so the standard $t$-structure gives
\[
  T^j(\mathcal G)
  =\Hom_{\D(\OO_P)}(\mathcal G,\omega_P[N+j])=0.
\]
On the other hand, $-j>N=\dim P$, so Grothendieck vanishing gives
$H^{-j}(P,\mathcal G)=0$; see
\cite[Chapter III, Theorem 2.7]{HartshorneAG}. Hence $S^j(\mathcal G)=0$ by
Lemma~\ref{lem:field-duality}. Thus $\eta_{\mathcal G}^j$ is an isomorphism
for every $j\in\mathbf Z$. Lemma~\ref{lem:devissage} now implies that
$\eta_G^j$ is an isomorphism for every $G\in\Dbcoh(P)$ and every
$j\in\mathbf Z$, proving \eqref{eq:projective-derived}.
\end{proof}

\begin{remark}[Adjunction and Neeman's calculation]
Grothendieck duality provides a right adjoint $p^\times$ to $\mathbf Rp_*$ on
the unbounded quasi-coherent derived category. The standard projective-space
calculation identifies its value at $k$ as
\[
  p^\times k\simeq\omega_P[N].
\]
Thus Theorem~\ref{thm:projective-derived} is the restriction to bounded
coherent complexes of the adjunction
\[
  \Hom_{\Dbcoh(P)}(G,p^\times k[j])
  \simeq\Hom_{\D(k)}(\mathbf Rp_*G,k[j]).
\]
This is the mechanism isolated in \cite[Section 3.3]{NeemanSimple}; the
geometric calculation of $p^\times k$ is supplied here by projective-space
duality.
\end{remark}

\section{The dualizing complex of a projective scheme}

Let $X$ be a projective scheme over $k$, and fix a closed immersion
\[
  i:X\hookrightarrow P=\PP^N_k.
\]

\begin{definition}[Exceptional inverse image for a closed immersion]
\label{def:closed-immersion-shriek}
For an $\OO_P$-module $M$, define
\[
  i^\flat M
  :=i^{-1}\shHom_{\OO_P}(i_*\OO_X,M).
\]
The $i_*\OO_X$-action on the sheaf Hom, given by precomposition, makes
$i^\flat M$ an $\OO_X$-module. The functor
\[
  i^\flat:\operatorname{Mod}(\OO_P)\longrightarrow
  \operatorname{Mod}(\OO_X)
\]
is left exact and right adjoint to the exact functor $i_*$. Its right-derived
functor is the exceptional inverse image
\[
  i^!:=\mathbf R i^\flat:
  \D(\OO_P)\longrightarrow\D(\OO_X).
\]
Thus, if $G\to J^\bullet$ is a K-injective resolution, then
\[
  i^!G\simeq i^\flat J^\bullet.
\]
Equivalently, after applying the fully faithful functor $i_*$, there is a
canonical isomorphism
\[
  i_*i^!G\simeq\RshHom_P(i_*\OO_X,G).
\]
If $G\in\D^+_{\mathrm{qc}}(P)$, then
$i^!G\in\D^+_{\mathrm{qc}}(X)$. More concretely, let
$U=\Spec R\subset P$ and $X\cap U=\Spec A$, where $A=R/I$. Under the affine
equivalences for quasi-coherent complexes, $(i^!G)|_{\Spec A}$ corresponds to
\[
  \RHom_R(A,G_U)\in\D^+(A),
\]
where $G_U\in\D^+(R)$ corresponds to $G|_U$ and the $A$-module structure on
the derived Hom is induced by the first argument.
\end{definition}

This is the closed-immersion case of Grothendieck duality; see
\cite[Tags 0A76, 0A78, and 0A9X]{Stacks}, and
\cite[Equation (4.9.4.1)]{LipmanNotes}.

\begin{proposition}[Derived adjunction for a closed immersion]
\label{prop:closed-adjunction}
For $F\in\Dbcoh(X)$ and $G\in\D^+_{\mathrm{qc}}(P)$ there are functorial
isomorphisms
\begin{equation}\label{eq:closed-internal-adjunction}
  i_*\RshHom_X(F,i^!G)
  \simeq\RshHom_P(i_*F,G)
\end{equation}
and
\begin{equation}\label{eq:closed-adjunction}
  \RHom_X(F,i^!G)\simeq\RHom_P(i_*F,G).
\end{equation}
In particular, on this range $i^!$ is right adjoint to the exact functor
$i_*$.
\end{proposition}

\begin{proof}
By Definition~\ref{def:closed-immersion-shriek}, the underived functor
$i^\flat$ is right adjoint to $i_*$. Since $i_*$ is exact for a closed
immersion, deriving this adjunction gives $\mathbf Ri_*=i_*$ and
\[
  \Hom_{\D(\OO_X)}(F,i^!G)
  \simeq\Hom_{\D(\OO_P)}(i_*F,G).
\]
The same argument at the level of Hom complexes proves the stronger
internal isomorphism \eqref{eq:closed-internal-adjunction}: if $J$ is a K-injective
representative of $G$, then $i^\flat J$ is K-injective on $X$, and the
ordinary adjunction gives an isomorphism of sheaf Hom complexes
\[
  i_*\shHom_X^\bullet\bigl(F,i^\flat J\bigr)
  \simeq \shHom_P^\bullet(i_*F,J).
\]
Applying derived global sections gives \eqref{eq:closed-adjunction}. See
\cite[Tags 0A75, 0A76, and 0E2I]{Stacks} and
\cite[Corollary 4.4.2]{LipmanNotes}.

For comparison, on an affine open the assertion is the following derived
tensor--Hom adjunction. If $A=R/I$, $M\in\D(A)$, and $J\in\D(R)$, then
\[
  \RHom_A\bigl(M,\RHom_R(A,J)\bigr)
  \simeq\RHom_R(M,J).
\]
Here the right-hand side regards $M$ as an $R$-complex. Finally,
\cite[Tag 0A9X]{Stacks} shows that this adjunction restricts to the stated
bounded-below quasi-coherent range.
\end{proof}

\begin{definition}[Projective dualizing complex]
\label{def:dualizing-complex}
The dualizing complex attached to the embedding $i$ is
\begin{equation}\label{eq:def-dualizing-complex}
  \omega_X^\bullet
  :=i^!\bigl(\omega_P[N]\bigr)
  \in\D(\OO_X).
\end{equation}
Equivalently,
\[
  i_*\omega_X^\bullet
  \simeq\RshHom_P(i_*\OO_X,\omega_P)[N].
\]
\end{definition}

To justify the asserted boundedness and coherence, note that $P$ is regular and
$i_*\OO_X$ has bounded coherent cohomology, hence is perfect on $P$; see
\cite[Tag 0FDC]{Stacks}. Therefore
$\RshHom_P(i_*\OO_X,\omega_P)$ has bounded coherent cohomology, and the same
is true after viewing it on $X$: exact pushforward along a closed immersion
identifies coherent sheaves on $X$ with coherent sheaves on $P$ annihilated
by the ideal of $X$; see \cite[Tag 01XY]{Stacks}.

\begin{theorem}[Derived duality for a projective scheme]
\label{thm:derived-duality-projective}
For every $F\in\Dbcoh(X)$ and every $j\in\mathbf Z$, there is a functorial
isomorphism
\begin{equation}\label{eq:derived-duality-projective}
  \Hom_{\Dbcoh(X)}(F,\omega_X^\bullet[j])
  \simeq
  \Hom_{\D(k)}\bigl(\RGamma(X,F),k[j]\bigr).
\end{equation}
In particular, for a coherent sheaf $\FF$,
\[
  \Hom_{\Dbcoh(X)}(\FF,\omega_X^\bullet[j])
  \simeq H^{-j}(X,\FF)\dual.
\]
If two projective embeddings are chosen, the resulting dualizing complexes
are canonically isomorphic: there is a unique isomorphism compatible with
the displayed duality isomorphisms.
\end{theorem}

\begin{proof}
By Definition \ref{def:dualizing-complex}, the fact that $i^!$ commutes with
shifts, Definition \ref{def:derived-hom}, and Proposition
\ref{prop:closed-adjunction}, we have
\begin{align*}
  \Hom_{\Dbcoh(X)}(F,\omega_X^\bullet[j])
  &=\Hom_{\D(\OO_X)}
      \bigl(F,i^!(\omega_P[N])[j]\bigr)\\
  &\simeq\Hom_{\D(\OO_X)}
      \bigl(F,i^!(\omega_P[N+j])\bigr)\\
  &=H^0\RHom_X\bigl(F,i^!(\omega_P[N+j])\bigr)\\
  &\simeq H^0\RHom_P\bigl(i_*F,\omega_P[N+j]\bigr)\\
  &=\Hom_{\D(\OO_P)}\bigl(i_*F,\omega_P[N+j]\bigr)\\
  &=\Hom_{\Dbcoh(P)}\bigl(i_*F,\omega_P[N+j]\bigr)\\
  &\simeq\Hom_{\D(k)}\bigl(\RGamma(P,i_*F),k[j]\bigr).
\end{align*}
A closed immersion has exact pushforward, and
$\Gamma(P,i_*(-))=\Gamma(X,-)$. Deriving this equality gives
$\RGamma(P,i_*F)\simeq\RGamma(X,F)$. Consequently,
\[
  \Hom_{\D(k)}\bigl(\RGamma(P,i_*F),k[j]\bigr)
  \simeq
  \Hom_{\D(k)}\bigl(\RGamma(X,F),k[j]\bigr).
\]
See
\cite[Chapter III, Exercise 4.1]{HartshorneAG} or
\cite[Tag 01XC]{Stacks}. This proves
\eqref{eq:derived-duality-projective}.

If $F=\FF$ is a coherent sheaf, then $\RGamma(X,\FF)$ is a bounded complex
with finite-dimensional cohomology. Lemma \ref{lem:field-duality} therefore
gives the stated identification with $H^{-j}(X,\FF)\dual$.

If two projective embeddings produce complexes $D$ and $D'$, the preceding
formula with $j=0$ identifies both as representing objects for the functor
\[
  \Phi(F)=\Hom_{\D(k)}\bigl(\RGamma(X,F),k\bigr)
  \qquad(F\in\Dbcoh(X)).
\]
The Yoneda lemma gives a unique isomorphism $D\simeq D'$ compatible with
these representing isomorphisms. This proves precisely the asserted
embedding independence; the compatibility requirement is essential for
uniqueness.
\end{proof}

\begin{remark}[Intrinsic interpretation via exceptional inverse image]
Let $p:P\to\Spec k$ be the structure morphism and let
$f=p\circ i:X\to\Spec k$:
\[
\begin{tikzcd}[column sep=large]
  X \arrow[r,hook,"i"] \arrow[dr,"f"'] & P \arrow[d,"p"] \\
  & \Spec k.
\end{tikzcd}
\]
We explain separately the intrinsic meaning of $\omega_X^\bullet$ and its
description by internal derived Hom.

For a proper morphism $g:Y\to Z$ in the range considered here, Grothendieck
duality provides an exceptional inverse image $g^!$, which agrees with the
right adjoint $g^\times$ of derived pushforward. Thus it is characterized by
functorial isomorphisms
\[
  \Hom_{\D(\OO_Y)}(A,g^!B)
  \simeq
  \Hom_{\D(\OO_Z)}(\mathbf Rg_*A,B).
\]
For the existence of this right adjoint, see \cite[Tag 0A9E]{Stacks}.
For the structure morphism $p:P\to\Spec k$, the projective-space calculation
of the preceding section is
\[
  p^!k\simeq\omega_P[N].
\]
Equivalently, $\omega_P[N]$ represents the restriction to $\Dbcoh(P)$ of the
functor
\[
  G\longmapsto
  \Hom_{\D(k)}\bigl(\RGamma(P,G),k\bigr).
\]

The intrinsic dualizing complex of $X$ is, by definition, $f^!k$; this
definition does not mention a projective embedding. Transitivity of
exceptional inverse image for the composite $f=p\circ i$ gives a natural
isomorphism of functors \cite[Tag 0ATX]{Stacks}:
\[
  (p\circ i)^!\simeq i^!\circ p^!.
\]
Evaluating it at $k$ and then using projective-space duality gives
\begin{align*}
  f^!k
  &=(p\circ i)^!k\\
  &\simeq i^!(p^!k)\\
  &\simeq i^!\bigl(\omega_P[N]\bigr)\\
  &=\omega_X^\bullet.
\end{align*}
Thus the embedding-dependent expression in Definition
\ref{def:dualizing-complex} is canonically isomorphic to the intrinsic object
$f^!k$.

We next derive the internal-Hom formula. Apply the internal adjunction
\eqref{eq:closed-internal-adjunction} with $F=\OO_X$ and
$G=\omega_P[N]$. It gives
\begin{align*}
  i_*\omega_X^\bullet
  &=i_*\RshHom_X\bigl(\OO_X,i^!(\omega_P[N])\bigr)\\
  &\simeq\RshHom_P(i_*\OO_X,\omega_P[N])\\
  &\simeq\RshHom_P(i_*\OO_X,\omega_P)[N].
\end{align*}
Here the first equality uses
$\RshHom_X(\OO_X,M)\simeq M$, while the last isomorphism uses the
compatibility of derived Hom with shifts in its second variable. Hence
\begin{equation}\label{eq:intrinsic-dualizing-complex}
  \boxed{
  i_*\omega_X^\bullet
  \simeq
  \RshHom_P(i_*\OO_X,\omega_P)[N].
  }
\end{equation}
This formula is therefore not an independent computation: it is the result
of projective-space duality, transitivity of $(-)^!$, and the closed-immersion
adjunction $i_*\dashv i^!$.

Under the canonical identification $\omega_X^\bullet\simeq f^!k$, derived
duality is precisely the adjunction
\begin{align*}
  \Hom_{\Dbcoh(X)}(F,\omega_X^\bullet[j])
  &\simeq
    \Hom_{\Dbcoh(X)}(F,f^!k[j])\\
  &\simeq
    \Hom_{\D(k)}(\mathbf Rf_*F,k[j])\\
  &=
    \Hom_{\D(k)}(\RGamma(X,F),k[j]).
\end{align*}
Here we used $\mathbf Rf_*F=\RGamma(X,F)$ and the compatibility of $f^!$
with shifts. In the notation $f^\times$ for the right adjoint of
$\mathbf Rf_*$, one has $f^!=f^\times$ on the present range.
If $X$ is smooth and proper of dimension $d$, the geometric identification
$f^\times k\simeq\Omega^d_{X/k}[d]$ recovers the usual smooth form of Serre
duality; see \cite[Tag 0FVV]{Stacks}. For the right-adjoint formulation, see
\cite[Section 3.3]{NeemanSimple}.
\end{remark}

\begin{definition}[Trace]
\label{def:trace}
Under the canonical identification $\omega_X^\bullet\simeq f^!k$, define the
derived trace morphism to be the counit of the adjunction
$\mathbf Rf_*\dashv f^!$, evaluated at $k$:
\[
  \operatorname{Tr}_X:
  \RGamma(X,\omega_X^\bullet)
  \simeq \mathbf Rf_*f^!k
  \longrightarrow k.
\]
Under the adjunction isomorphism
\[
  \Hom_{\D(k)}\bigl(\mathbf Rf_*f^!k,k\bigr)
  \simeq
  \Hom_{\D(\OO_X)}(f^!k,f^!k),
\]
the morphism $\operatorname{Tr}_X$ corresponds to
$\mathrm{id}_{f^!k}$. Thus this definition agrees with the morphism
corresponding to $\mathrm{id}_{\omega_X^\bullet}$ under
\eqref{eq:derived-duality-projective}.

If $X$ is Cohen--Macaulay of pure dimension $d$, Theorem
\ref{thm:CM-concentration} below will give
$\omega_X^\bullet\simeq\omega_X[d]$. Under this identification, taking
degree-zero cohomology gives
\[
  H^0\RGamma(X,\omega_X^\bullet)
  \simeq H^0\RGamma(X,\omega_X[d])
  =H^d(X,\omega_X),
\]
so $H^0(\operatorname{Tr}_X)$ is the classical trace map
\[
  \tr_X:H^d(X,\omega_X)\longrightarrow k.
\]
See
\cite[Tag 0A9D]{Stacks} and
\cite[Section 3.3]{NeemanSimple}.
\end{definition}

\section{Cohen--Macaulay local algebra}

The concentration of $\omega_X^\bullet$ is local.  The required commutative
algebra is recalled in a form adapted to the embedding
$X\hookrightarrow\PP^N_k$.

\begin{definition}[Regular sequence and depth]
Let $(R,\mathfrak m)$ be a noetherian local ring and let $M$ be a finite
$R$-module.

A sequence $x_1,\ldots,x_r\in\mathfrak m$ is \emph{$M$-regular} if
$x_j$ is a non-zero-divisor on
$M/(x_1,\ldots,x_{j-1})M$ for every $j$, and
$M/(x_1,\ldots,x_r)M\neq 0$.
The depth $\depth_R M$ is the maximal length of an $M$-regular sequence
contained in $\mathfrak m$.
\end{definition}

\begin{definition}[Cohen--Macaulay ring and scheme]
A noetherian local ring $(R,\mathfrak m)$ is \emph{Cohen--Macaulay} if
\[
  \depth R=\dim R.
\]
A locally noetherian scheme $X$ is Cohen--Macaulay if every local ring
$\OO_{X,x}$ is Cohen--Macaulay.  It is \emph{equidimensional of dimension
$d$} if every irreducible component has dimension $d$.
\end{definition}

See \cite[Chapter III, Section 6]{HartshorneAG},
\cite[Chapter 2]{BrunsHerzog}, and
\cite[Section 17]{Matsumura} for these notions and their basic properties.

\begin{definition}[Grade]
For a proper ideal $I\subset R$, its grade is
\[
  \grade(I,R)
  :=\inf\{q\geq 0:\Ext_R^q(R/I,R)\neq 0\}.
\]
Equivalently, it is the maximal length of an $R$-regular sequence contained
in $I$.
\end{definition}

The equivalence of the two descriptions is proved in
\cite[Theorem 1.2.5]{BrunsHerzog}.  If $R$ is Cohen--Macaulay, then
\begin{equation}\label{eq:grade-height}
  \grade(I,R)=\height(I)
\end{equation}
for every proper ideal $I$; see
\cite[Theorem 17.4]{Matsumura}.

\begin{theorem}[Auslander--Buchsbaum]
\label{thm:AB}
Let $(R,\mathfrak m)$ be a noetherian local ring and let $M$ be a nonzero
finite $R$-module of finite projective dimension.  Then
\[
  \pd_R(M)+\depth_R(M)=\depth R.
\]
\end{theorem}

See \cite[Chapter III, Proposition 6.12A]{HartshorneAG},
\cite[Theorem 1.3.3]{BrunsHerzog}, or
\cite[Theorem 19.1]{Matsumura}.

\begin{proposition}[Ext concentration in a regular ambient ring]
\label{prop:local-ext-concentration}
Let $(R,\mathfrak m)$ be a regular local ring of dimension $n$, let
$I\subsetneq R$, and put $A=R/I$.  Assume that $A$ is Cohen--Macaulay of
dimension $e$.  Set $c=n-e=\height I$.  Then
\begin{equation}\label{eq:local-ext-concentration}
  \Ext_R^q(A,R)=0\qquad(q\neq c).
\end{equation}
Moreover,
\[
  \omega_A:=\Ext_R^c(A,R)
\]
is a maximal Cohen--Macaulay $A$-module with
$\Supp_A(\omega_A)=\Spec A$.
\end{proposition}

\begin{proof}
A regular local ring has finite global dimension $n$ and is
Cohen--Macaulay; see
\cite[Chapter III, Propositions 6.11A and 6.12A]{HartshorneAG} and
\cite[Theorems 17.8 and 19.2]{Matsumura}. Hence $A$ has finite projective dimension
over $R$.

Since $R$ is Cohen--Macaulay, $\dim(R/I)=n-\height I$; see
\cite[Tags 00N9 and 00NA]{Stacks}. Hence $c=n-e=\height I$, as asserted.

The maximal ideal of $A$ is $\mathfrak mA$, so
\[
  \depth_R A=\depth_A A=e
\]
because a sequence in $\mathfrak m$ is $A$-regular exactly when its image
in $\mathfrak mA$ is $A$-regular, and $A$ is Cohen--Macaulay. Applying
Theorem \ref{thm:AB} gives
\[
  \pd_R A=n-e=c.
\]
Therefore $\Ext_R^q(A,R)=0$ for $q>c$.

For the lower degrees, the definition of grade and
\eqref{eq:grade-height} give
\[
  \inf\{q:\Ext_R^q(A,R)\neq 0\}
  =\grade(I,R)=\height I=c.
\]
Thus $\Ext_R^q(A,R)=0$ for $q<c$, proving
\eqref{eq:local-ext-concentration}.

It remains to prove the assertions about the surviving Ext module. Because
$R$ is regular local, it is Gorenstein and $R[n]$ is a normalized dualizing
complex for $R$. For the finite surjection $R\to A$, finite-map duality gives
the normalized dualizing complex
\[
  \omega_A^\bullet
  :=\RHom_R\bigl(A,R[n]\bigr)
  \simeq \RHom_R(A,R)[n]
\]
for $A$; see \cite[Tag 0A7M]{Stacks}. Its cohomology is
\[
  H^r(\omega_A^\bullet)
  =\Ext_R^{n+r}(A,R).
\]
By \eqref{eq:local-ext-concentration}, this vanishes unless
$n+r=c=n-e$, or equivalently $r=-e$, and in that degree it equals
$\omega_A=\Ext_R^c(A,R)$. Hence
\[
  \omega_A^\bullet\simeq\omega_A[e].
\]
Since $A$ is Cohen--Macaulay, its dualizing module $\omega_A$ is maximal
Cohen--Macaulay; see \cite[Tag 0AWS]{Stacks}. Moreover, the support of a dualizing module is the union of
the irreducible components of maximal dimension. A Cohen--Macaulay local ring
is equidimensional, so every irreducible component of $\Spec A$ has dimension
$e$; consequently $\Supp_A(\omega_A)=\Spec A$; see
\cite[Tag 0AWH]{Stacks}.
\end{proof}

\begin{remark}
The two halves of the proof have different origins.  Auslander--Buchsbaum
controls the upper vanishing $q>c$, whereas grade equals height controls the
lower vanishing $q<c$.  Both are needed to obtain concentration in exactly
one degree.
\end{remark}

\section{Concentration of the dualizing complex}

Let $X\hookrightarrow P=\PP^N_k$ be a projective, equidimensional
Cohen--Macaulay scheme of dimension $d$, and set $c=N-d$.

\begin{theorem}[Cohen--Macaulay concentration]
\label{thm:CM-concentration}
For $q\geq0$, the sheaves
\[
  \shExt_P^q(i_*\OO_X,\omega_P)
\]
vanish for $q\neq c$. Define the canonical sheaf intrinsically by
\begin{equation}\label{eq:def-canonical-sheaf}
  \omega_X:=H^{-d}(\omega_X^\bullet).
\end{equation}
Then
\begin{equation}\label{eq:CM-concentration}
  \omega_X^\bullet\simeq\omega_X[d].
\end{equation}
The sheaf $\omega_X$ is coherent, Cohen--Macaulay, and has support $X$.
\end{theorem}

\begin{proof}
The sheaves $\shExt_P^q(i_*\OO_X,\omega_P)$ are coherent, and their stalks
are computed by local Ext; see
\cite[Chapter III, Proposition 6.8]{HartshorneAG}. Because $P$ is universally
catenary and $X$ is equidimensional of dimension $d$, the closed subscheme
$X\subset P$ has pure codimension $c=N-d$; see
\cite[Tags 02JB and 02I6]{Stacks}.

Now let $x\in X$ be arbitrary and put
\[
  R=\OO_{P,x},\qquad A=\OO_{X,x}=R/I_x.
\]
The ring $R$ is regular, $A$ is Cohen--Macaulay, and the pure-codimension
statement gives
\[
  \dim R-\dim A=\height(I_x)=c.
\]
Using the stalk formula above and the invertibility of $\omega_P$, we obtain
\[
  \shExt_P^q(i_*\OO_X,\omega_P)_x
  \simeq\Ext_R^q\bigl(A,(\omega_P)_x\bigr)
  \simeq\Ext_R^q(A,R)\otimes_R(\omega_P)_x.
\]
Proposition \ref{prop:local-ext-concentration} gives vanishing for $q\neq c$.

By Definition \ref{def:dualizing-complex} and the exactness of $i_*$,
\[
  i_*H^r(\omega_X^\bullet)
  \simeq H^r(i_*\omega_X^\bullet)
  \simeq\shExt_P^{N+r}(i_*\OO_X,\omega_P).
\]
The only nonzero term occurs when $N+r=c=N-d$, i.e. when $r=-d$. Thus
$H^r(\omega_X^\bullet)=0$ for $r\neq-d$, while
$H^{-d}(\omega_X^\bullet)=\omega_X$ by
\eqref{eq:def-canonical-sheaf}. The canonical truncation morphism therefore
yields
$\omega_X^\bullet\simeq\omega_X[d]$.

Combining the cohomology identity above in degree $r=-d$ with the displayed
stalk calculation identifies $(\omega_X)_x$ with the canonical module
$\Ext_R^c(A,R)$ tensored by the free rank-one module $(\omega_P)_x$.
Proposition
\ref{prop:local-ext-concentration} shows that this stalk is maximal
Cohen--Macaulay and has support $\Spec A$. Since $x\in X$ was arbitrary,
$\omega_X$ is coherent and Cohen--Macaulay and its support is all of $X$.
See also \cite[Tags 0AWT and 0FVZ]{Stacks}.
\end{proof}

\begin{corollary}[Canonical sheaf as an Ext sheaf]
\label{cor:canonical-sheaf-ext}
Under the closed immersion $i:X\hookrightarrow P=\PP_k^N$, the canonical
sheaf of the Cohen--Macaulay scheme $X$ satisfies
\[
  i_*\omega_X
  \simeq
  \shExt_P^{N-d}(i_*\OO_X,\omega_P).
\]
Thus the canonical sheaf is the Ext sheaf in the codimension
$N-d$ of $X$ in $P$.
\end{corollary}

\begin{proof}
Theorem \ref{thm:CM-concentration} gives
$\omega_X^\bullet\simeq\omega_X[d]$, so
\[
  H^{-d}(\omega_X^\bullet)=\omega_X.
\]
Because $i_*$ is exact, taking degree-$(-d)$ cohomology in
\eqref{eq:intrinsic-dualizing-complex} gives $i_*\omega_X$ on the left. On
the right, the shift convention yields
\begin{align*}
  H^{-d}\!\left(
    \RshHom_P(i_*\OO_X,\omega_P)[N]
  \right)
  &=H^{N-d}\RshHom_P(i_*\OO_X,\omega_P)\\
  &=\shExt_P^{N-d}(i_*\OO_X,\omega_P),
\end{align*}
which proves the assertion.
\end{proof}

\begin{corollary}[Identification with the dualizing sheaf]
\label{cor:hartshorne-dualizing-sheaf}
For every coherent sheaf $\FF$ on $X$, there is a functorial isomorphism
\begin{equation}\label{eq:top-duality}
  \Hom_X(\FF,\omega_X)
  \simeq H^d(X,\FF)\dual.
\end{equation}
Hence $\omega_X$ is a dualizing sheaf in the sense of
\cite[Chapter III, Definition preceding Proposition 7.2]{HartshorneAG}.
In particular, it is uniquely isomorphic to any dualizing sheaf obtained by
the finite-morphism construction, provided the isomorphism is required to be
compatible with the two representing isomorphisms.
\end{corollary}

\begin{proof}
Set $j=-d$ in Theorem \ref{thm:derived-duality-projective}.  Using
$\omega_X^\bullet\simeq\omega_X[d]$ gives
\[
  \Hom_{\Dbcoh(X)}(\FF,\omega_X)
  \simeq H^d(X,\FF)\dual.
\]
Since both arguments are sheaves placed in degree zero, the left-hand side is
$\Hom_X(\FF,\omega_X)$.  Uniqueness follows from the Yoneda lemma.
\end{proof}

\begin{remark}[Intrinsic meaning of Cohen--Macaulayness]
For a locally noetherian scheme carrying a dualizing complex,
Cohen--Macaulayness is equivalent to the local concentration of that complex
in a single degree; see \cite[Tag 0AWT]{Stacks}.  Theorem
\ref{thm:CM-concentration} is the projective, equidimensional form of this
criterion, with the unique nonzero degree fixed to be $-d$.
\end{remark}

\section{Serre duality and the canonical pairing}
\label{sec:serre-duality-pairing}

\begin{theorem}[Serre duality for Cohen--Macaulay schemes]
\label{thm:main-serre}
Let $X$ be a projective, equidimensional Cohen--Macaulay scheme of dimension
$d$ over a field $k$, and let $\omega_X$ be a dualizing sheaf on $X$. Then,
for every coherent sheaf $\FF$ and every $i\geq0$, there is a functorial
isomorphism
\begin{equation}\label{eq:main-serre}
  \theta^i_{\FF}:\Ext_X^i(\FF,\omega_X)
  \xrightarrow{\ \sim\ }H^{d-i}(X,\FF)\dual.
\end{equation}
Here $H^q(X,\FF)$ is understood to be zero for $q<0$. For
$0\leq i\leq d$, the isomorphism is induced by the Yoneda--trace pairing
\begin{equation}\label{eq:main-pairing}
  \Ext_X^i(\FF,\omega_X)\times H^{d-i}(X,\FF)
  \longrightarrow H^d(X,\omega_X)\xrightarrow{\tr_X}k,
\end{equation}
and this pairing is perfect. For $i>d$, both vector spaces in
\eqref{eq:main-serre} vanish.
\end{theorem}

This is the implication $(i)\Rightarrow(iii)$ in
\cite[Chapter III, Theorem 7.6(b)]{HartshorneAG}. Hartshorne assumes that
$k$ is algebraically closed; the argument below works over an arbitrary
field. A formulation in terms of dualizing complexes appears in
\cite[Tag 0FVZ]{Stacks}.

\begin{proof}[Proof of Theorem \ref{thm:main-serre}]
Let
\[
  \omega_X^{\mathrm{can}}:=H^{-d}(\omega_X^\bullet).
\]
Theorem \ref{thm:CM-concentration} and Corollary
\ref{cor:hartshorne-dualizing-sheaf} show that
$\omega_X^{\mathrm{can}}$ is a dualizing sheaf. Since dualizing sheaves are
unique up to a unique isomorphism compatible with their representing
properties, there is a canonical isomorphism
$\omega_X^{\mathrm{can}}\simeq\omega_X$. We use this identification below.

For a coherent sheaf $\FF$ and $i\geq0$, Proposition
\ref{prop:ext-derived}, Theorem \ref{thm:CM-concentration}, and Theorem
\ref{thm:derived-duality-projective} give
\begin{align*}
  \Ext_X^i(\FF,\omega_X)
  &\simeq\Hom_{\Dbcoh(X)}(\FF,\omega_X[i])\\
  &\simeq\Hom_{\Dbcoh(X)}(\FF,\omega_X^\bullet[i-d])\\
  &\simeq H^{d-i}(X,\FF)\dual.
\end{align*}
The construction is functorial in $\FF$.

Assume first that $0\leq i\leq d$. Let
\[
  \alpha\in\Ext_X^i(\FF,\omega_X),
  \qquad
  \beta\in H^{d-i}(X,\FF)
  \simeq\Ext_X^{d-i}(\OO_X,\FF).
\]
Using Proposition \ref{prop:ext-derived}, represent these classes by
morphisms in the derived category
\[
  \beta:\OO_X\longrightarrow\FF[d-i],
  \qquad
  \alpha:\FF\longrightarrow\omega_X[i].
\]
After shifting $\alpha$ by $d-i$, their composite is
\[
  \OO_X\xrightarrow{\ \beta\ }\FF[d-i]
  \xrightarrow{\ \alpha[d-i]\ }\omega_X[d].
\]
It represents their Yoneda product
\[
  \alpha\circ\beta
  \in\Ext_X^d(\OO_X,\omega_X)=H^d(X,\omega_X).
\]
The duality isomorphism is induced by the adjunction counit, which is the
trace map of Definition \ref{def:trace}. Naturality of adjunction therefore sends
$\alpha$ precisely to the functional
\[
  \beta\longmapsto\tr_X(\alpha\circ\beta).
\]
This identifies \eqref{eq:main-serre} with the pairing
\eqref{eq:main-pairing}.

If $i>d$, then $d-i<0$, so $H^{d-i}(X,\FF)=0$. The duality isomorphism
already proved in \eqref{eq:main-serre} therefore also gives
$\Ext_X^i(\FF,\omega_X)=0$. This explains why no negative ordinary Ext group
is needed in the argument.

For $0\leq i\leq d$, the group $H^{d-i}(X,\FF)$ is finite-dimensional
because $X$ is projective and $\FF$ is coherent
\cite[Chapter III, Theorem 5.2]{HartshorneAG}; for $i>d$ it is zero. Hence
\eqref{eq:main-serre} also makes $\Ext_X^i(\FF,\omega_X)$ finite-dimensional.
An isomorphism from the first space to the dual of the second therefore
induces an isomorphism from the second space to the dual of the first. The
pairing is perfect in both variables.
\end{proof}

\begin{corollary}[Classical vector-bundle form]
\label{cor:vector-bundle-duality}
Let $\EE$ be a vector bundle on $X$. For every $0\leq i\leq d$, there is a
functorial isomorphism
\[
  H^i(X,\EE\dual\otimes\omega_X)
  \simeq H^{d-i}(X,\EE)\dual.
\]
Equivalently, the cup-product and trace define a perfect pairing
\[
  H^i(X,\EE\dual\otimes\omega_X)
  \times H^{d-i}(X,\EE)
  \longrightarrow H^d(X,\omega_X)
  \xrightarrow{\tr_X}k.
\]
\end{corollary}

\begin{proof}
Since $\EE$ is locally free,
\[
  \RshHom_X(\EE,\omega_X)
  \simeq\EE\dual\otimes\omega_X.
\]
Consequently
\[
  \Ext_X^i(\EE,\omega_X)
  \simeq H^i(X,\EE\dual\otimes\omega_X).
\]
Apply Theorem \ref{thm:main-serre}.  Compare
\cite[Chapter III, Corollary 7.7]{HartshorneAG}.
\end{proof}

\begin{remark}[Agreement with Hartshorne's maps]
The calculation in the proof shows directly that $\theta^i_{\FF}$ sends a
class $\alpha$ to the functional
\[
  \beta\longmapsto\tr_X(\alpha\circ\beta).
\]
This is precisely the natural transformation constructed from the trace
pairing in \cite[Chapter III, Theorem 7.6]{HartshorneAG}; hence the derived
isomorphisms agree with Hartshorne's maps.
\end{remark}

\begin{remark}[Classical and derived proofs]
Hartshorne derives the theorem by two steps: the Cohen--Macaulay condition
produces cohomological vanishing for sufficiently negative twists, and that
vanishing proves effaceability of the relevant $\delta$-functor.  The derived
proof packages the same local algebra into the statement
$\omega_X^\bullet\simeq\omega_X[d]$.  Once this concentration is known, all
degrees of duality follow simultaneously from one adjunction. See
\cite[Chapter III, proof of Theorem 7.6]{HartshorneAG}.
\end{remark}

\section{Local complete intersections}

The Cohen--Macaulay hypothesis determines the amplitude of the dualizing
complex but does not usually identify the sheaf $\omega_X$ explicitly.  A
local complete intersection admits an explicit calculation because its
structure sheaf has a Koszul resolution.

\begin{definition}[Regular immersion and local complete intersection]
Let $i:X\hookrightarrow Y$ be a closed immersion of locally noetherian
schemes, with ideal sheaf $\II$. The immersion is \emph{regular of pure
codimension $c$} if, locally on $Y$ near every point of $X$, the ideal $\II$
is generated by a regular sequence of length $c$. In this case
$\II/\II^2$ is locally free of rank $c$; it is the conormal bundle, and
\[
  N_{X/Y}:=(\II/\II^2)\dual
\]
is the normal bundle.

A locally noetherian $k$-scheme $X$ is a \emph{local complete intersection
over $k$} if every point has an open neighbourhood that admits a regular
closed immersion into a smooth $k$-scheme.
\end{definition}

See \cite[Chapter II, Section 8]{HartshorneAG} and
\cite[Tags 0638 and 067M]{Stacks}.

\begin{proposition}[Koszul resolution]
\label{prop:koszul-resolution}
Let $R$ be a ring and let $f_1,\ldots,f_c$ be an $R$-regular sequence.
Put $I=(f_1,\ldots,f_c)$, $A=R/I$, and let $F=R^c$ with basis
$e_1,\ldots,e_c$. Give the Koszul complex the differential determined by
$e_j\mapsto f_j$. Then
\[
  K(f_1,\ldots,f_c):
  0\to\bigwedge^cF\to\cdots\to\bigwedge^1F\to R\to0
\]
is a finite free resolution of $A$. Consequently, for $q\geq0$,
\begin{equation}\label{eq:koszul-ext}
  \Ext_R^q(A,R)=0\quad(q\neq c),
  \qquad
  \Ext_R^c(A,R)
  \simeq\determinant_A(I/I^2)\dual.
\end{equation}
\end{proposition}

\begin{proof}
Exactness of the Koszul complex for a regular sequence is proved in
\cite[Chapter III, Proposition 7.10A]{HartshorneAG}; see also
\cite[Section 1.6]{BrunsHerzog}. Applying $\Hom_R(-,R)$ to the resolution
gives the dual Koszul complex, which is exact away from its final
cohomology. Its only cohomology is in degree $c$, where
it is
\[
  A\otimes_R\bigwedge^cF\dual.
\]
The canonical map
\[
  A\otimes_R F\longrightarrow I/I^2,
  \qquad e_j\longmapsto f_j\bmod I^2,
\]
is an isomorphism because $f_1,\ldots,f_c$ is a regular sequence. Taking its
top exterior power and dualizing yields \eqref{eq:koszul-ext}.
\end{proof}

\begin{remark}[Canonicity of the determinant factor]
The isomorphism in \eqref{eq:koszul-ext} is independent of the chosen regular
sequence. Indeed, two local regular generating sequences determine two bases
of the locally free $A$-module $I/I^2$. On an overlap these bases differ by a
matrix $M\in\operatorname{GL}_c(A)$. The top exterior generators differ by
$\det(M)$ and their duals by $\det(M)^{-1}$, exactly as prescribed by the
transition functions of $\determinant_A(I/I^2)\dual$. Hence the local
isomorphisms are canonical and glue.
\end{remark}

\begin{theorem}[Adjunction formula for a local complete intersection]
\label{thm:lci-adjunction}
Let $i:X\hookrightarrow P=\PP^N_k$ be a regular closed immersion of pure
codimension $c$, with ideal sheaf $\II$. Then $X$ is Cohen--Macaulay of pure
dimension $d=N-c$, and
\begin{equation}\label{eq:lci-adjunction}
  \omega_X
  \simeq
  \omega_P|_X\otimes
  \determinant(\II/\II^2)\dual
  \simeq\omega_P|_X\otimes\determinant(N_{X/P}).
\end{equation}
Since $\omega_P\simeq\determinant(\Omega_{P/k})$, this is equivalently
\begin{equation}\label{eq:lci-cotangent-conormal}
  \omega_X
  \simeq
  \determinant(\Omega_{P/k}|_X)
  \otimes\determinant(\II/\II^2)\dual.
\end{equation}
Equivalently,
\begin{equation}\label{eq:lci-shriek}
  i^!\omega_P
  \simeq
  \omega_P|_X\otimes\determinant(N_{X/P})[-c].
\end{equation}
\end{theorem}

\begin{proof}
At every $x\in X$, the regular local ring $\OO_{P,x}$ surjects onto
$\OO_{X,x}$ with kernel generated by a regular sequence of length $c$. A
quotient of a regular local ring by a regular sequence is Cohen--Macaulay of
the expected dimension; see
\cite[Chapter II, Theorem 8.21A(a), (c), and (d)]{HartshorneAG} or
\cite[Theorem 2.1.3]{BrunsHerzog}. Hence $X$ is Cohen--Macaulay of pure
dimension $N-c$, and Theorem
\ref{thm:CM-concentration} applies.

Corollary \ref{cor:canonical-sheaf-ext}, Proposition
\ref{prop:koszul-resolution}, and the local freeness of $\omega_P$ give
\begin{align*}
  i_*\omega_X
  &\simeq\shExt_P^c(i_*\OO_X,\omega_P)\\
  &\simeq
  i_*\!\left(
    \omega_P|_X\otimes\determinant(\II/\II^2)\dual
  \right).
\end{align*}
The second isomorphism is computed locally by the Koszul resolution, and the
change-of-generators calculation above shows that these local isomorphisms
glue. Since pushforward along a closed immersion is fully faithful, this
proves \eqref{eq:lci-adjunction}.

For the derived formula, the closed-immersion adjunction gives
\begin{align*}
  i_*i^!\omega_P
  &\simeq\RshHom_P(i_*\OO_X,\omega_P)\\
  &\simeq
  i_*\!\left(
    \omega_P|_X\otimes\determinant(N_{X/P})
  \right)[-c].
\end{align*}
Indeed, the Koszul calculation shows that the derived Hom on the first line
has only one nonzero cohomology sheaf, in degree $c$. With our convention
$K[m]^r=K^{m+r}$, a complex whose sole cohomology sheaf is $E$ in degree $c$
is isomorphic to $E[-c]$. Full faithfulness of $i_*$ now gives
\eqref{eq:lci-shriek}.

The sheaf formula \eqref{eq:lci-adjunction} is
\cite[Chapter III, Theorem 7.11]{HartshorneAG}; for the exceptional
inverse-image formula \eqref{eq:lci-shriek}, see \cite[Tag 0BR0]{Stacks}.
\end{proof}

\begin{corollary}[Gorenstein property]
Every local complete intersection scheme over $k$ is Gorenstein. In the
projective situation of Theorem \ref{thm:lci-adjunction}, its dualizing sheaf
is therefore invertible.
\end{corollary}

\begin{proof}
At a point of a local complete intersection over $k$, choose a regular closed
immersion into a smooth $k$-scheme. The ambient local ring is regular, and the
local ring of $X$ is its quotient by a regular sequence. It is
Cohen--Macaulay, and Proposition \ref{prop:koszul-resolution} identifies its
canonical module with the determinant of the dual conormal module, which is
free of rank one. Hence the local ring is Gorenstein. In the projective
setting, Formula \eqref{eq:lci-adjunction} expresses $\omega_X$ as a tensor
product of invertible sheaves. See
\cite[Chapter III, Theorem 7.11]{HartshorneAG} and
\cite[Tag 0BVA]{Stacks}.
\end{proof}

\begin{example}[Hypersurfaces]
Let $X\subset\PP^N_k$ be a hypersurface defined by a homogeneous polynomial of
degree $m$. Then $X$ is an effective Cartier divisor, so it is a local
complete intersection of codimension one, and
\[
  \II/\II^2\simeq\OO_X(-m).
\]
Since $\omega_{\PP^N_k}\simeq\OO_{\PP^N_k}(-N-1)$, the adjunction formula
gives
\[
  \omega_X
  \simeq\OO_X(-N-1)\otimes\OO_X(m)
  \simeq\OO_X(m-N-1).
\]
For example, a plane curve of degree $m$ has
$\omega_X\simeq\OO_X(m-3)$.
\end{example}

\section{The smooth case}

Let $X$ be a smooth and projective $k$-scheme of pure dimension $d$. Choose a
closed immersion $i:X\hookrightarrow P=\PP^N_k$. Since $X$ and $P$ are
smooth over $k$, $i$ is a regular immersion of pure codimension $c=N-d$; see
\cite[Chapter II, Theorem 8.17]{HartshorneAG} and
\cite[Tag 0H1G]{Stacks}.
We write
\[
  \Omega^d_{X/k}:=\bigwedge^d\Omega_{X/k}.
\]

\begin{proposition}[Conormal sequence]
\label{prop:conormal-sequence}
There is a short exact sequence of vector bundles on $X$
\begin{equation}\label{eq:conormal-sequence}
  0\longrightarrow \II/\II^2
  \longrightarrow \Omega_{P/k}|_X
  \longrightarrow \Omega_{X/k}
  \longrightarrow 0.
\end{equation}
\end{proposition}

\begin{proof}
The conormal sequence is always right exact.  For a regular immersion into a
smooth scheme with smooth source, the left-hand map is injective and all
three terms are locally free.  See
\cite[Chapter II, Proposition 8.12 and Theorem 8.17]{HartshorneAG} or
\cite[Tag 06BB]{Stacks}.
\end{proof}

\begin{theorem}[Canonical sheaf of a smooth projective scheme]
\label{thm:smooth-canonical}
For a smooth projective $k$-scheme $X$ of pure dimension $d$,
\begin{equation}\label{eq:smooth-canonical}
  \omega_X\simeq\determinant(\Omega_{X/k})=\Omega^d_{X/k}.
\end{equation}
Consequently, for every vector bundle $\EE$ and every $0\leq i\leq d$, cup
product, evaluation, and trace define a perfect pairing
\[
  H^i(X,\EE\dual\otimes\Omega^d_{X/k})
  \times H^{d-i}(X,\EE)
  \longrightarrow H^d(X,\Omega^d_{X/k})
  \xrightarrow{\tr_X}k.
\]
\end{theorem}

\begin{proof}
Taking determinants in \eqref{eq:conormal-sequence} gives
\[
  \determinant(\Omega_{P/k}|_X)
  \simeq
  \determinant(\II/\II^2)
  \otimes\determinant(\Omega_{X/k}).
\]
Since $\determinant(\Omega_{P/k})\simeq\omega_P$, Theorem
\ref{thm:lci-adjunction} yields
\begin{align*}
  \omega_X
  &\simeq
  \omega_P|_X\otimes\determinant(\II/\II^2)\dual\\
  &\simeq\determinant(\Omega_{X/k}).
\end{align*}
The pairing is Corollary \ref{cor:vector-bundle-duality} under this
identification.  Compare
\cite[Chapter III, Corollary 7.12]{HartshorneAG} and
\cite[Tag 0FVV]{Stacks}.
\end{proof}


\begin{thebibliography}{99}

\bibitem{BrunsHerzog}
W.~Bruns and H.~J.~Herzog,
\emph{Cohen--Macaulay Rings}, revised edition,
Cambridge Studies in Advanced Mathematics 39,
Cambridge University Press, 1998;
\url{https://doi.org/10.1017/CBO9780511608681}.

\bibitem{HartshorneAG}
R.~Hartshorne,
\emph{Algebraic Geometry},
Graduate Texts in Mathematics 52,
Springer, 1977;
\url{https://doi.org/10.1007/978-1-4757-3849-0}.

\bibitem{LipmanNotes}
J.~Lipman,
\emph{Notes on Derived Functors and Grothendieck Duality},
in J.~Lipman and M.~Hashimoto,
\emph{Foundations of Grothendieck Duality for Diagrams of Schemes},
Lecture Notes in Mathematics 1960,
Springer, 2009, pp.~1--259;
\url{https://doi.org/10.1007/978-3-540-85420-3}.

\bibitem{Matsumura}
H.~Matsumura,
\emph{Commutative Ring Theory}, translated by M.~Reid, first paperback edition with corrections,
Cambridge Studies in Advanced Mathematics 8,
Cambridge University Press, 1989;
\url{https://doi.org/10.1017/CBO9781139171762}.

\bibitem{NeemanSimple}
A.~Neeman,
\emph{Grothendieck duality made simple},
in G.~Corti\~nas and C.~A.~Weibel (eds.),
\emph{$K$-theory in Algebra, Analysis and Topology},
Contemporary Mathematics 749,
American Mathematical Society, 2020, pp.~279--325;
\url{https://doi.org/10.1090/conm/749/15076};
\href{https://arxiv.org/abs/1806.03293}{arXiv:1806.03293}.

\bibitem{Stacks}
The Stacks Project Authors,
\emph{The Stacks Project},
\href{https://stacks.math.columbia.edu/}{stacks.math.columbia.edu}.
Individual results are cited in the text by their stable tags; accessed
17 July 2026.

\bibitem{Weibel}
C.~A.~Weibel,
\emph{An Introduction to Homological Algebra},
Cambridge Studies in Advanced Mathematics 38,
Cambridge University Press, 1994;
\url{https://doi.org/10.1017/CBO9781139644136}.

\end{thebibliography}

\end{document}
